\section{Categoricity}
\label{sec:qe-categoricity-02}

\begin{definition}[Categoricity]
Let \(T\) be a theory and let \(\kappa\) be an infinite cardinal.  The theory
\(T\) is \emph{\(\kappa\)-categorical} iff \(T\) has a model of cardinality
\(\kappa\) and any two models of \(T\) of cardinality \(\kappa\) are
isomorphic.
\end{definition}

\begin{warning}
Categoricity at one cardinal is not categoricity at every cardinal.  The
L\"owenheim--Skolem theorems imply that a first-order theory with an infinite
model has models in multiple infinite cardinalities, subject to the language
bound.  They do not preclude uniqueness up to isomorphism separately within
each cardinality.  For example, the equality-only theory asserting that there
are at least \(n\) elements for every \(n\geq1\) is
\(\kappa\)-categorical for every infinite cardinal \(\kappa\): any two pure
equality structures of size \(\kappa\) are isomorphic.  A transfer theorem
such as Morley's theorem requires additional hypotheses before categoricity at
one cardinal implies categoricity at others.
\end{warning}

\begin{definition}[Finite partial isomorphism]
A \emph{finite partial isomorphism} from \(\mathcal M\) to \(\mathcal N\) is a
finite injective map \(f:A\to N\), with \(A\subseteq M\), that preserves and
reflects all atomic formulas on tuples from \(A\).
\end{definition}

\begin{method}[Back-and-forth]
\label{method:back-and-forth}
For countable structures, build an isomorphism as the union of an increasing
sequence of finite partial isomorphisms.  At each ``forth'' stage extend the
domain to contain the next enumerated element of the first structure; at each
``back'' stage extend the range to contain the next enumerated element of the
second.  The name denotes this alternating proof method, not a new map or
relation.  It differs from the one-sided extension property in
\cref{def:homogeneity} because it enforces totality and surjectivity together.
The next theorem gives the exact hypotheses, and the dense-order proof then
uses them.
\end{method}

\begin{theorem}[Back-and-forth test]
\label{thm:back-and-forth-test}
Let \(\mathcal M\) and \(\mathcal N\) be countable structures.  Suppose
\(\mathcal F\) is a nonempty set of finite partial isomorphisms from
\(\mathcal M\) to \(\mathcal N\) such that, for every \(f\in\mathcal F\):
\begin{enumerate}
  \item for every \(a\in M\), some \(g\in\mathcal F\) extends \(f\) and has
        \(a\) in its domain; and
  \item for every \(b\in N\), some \(g\in\mathcal F\) extends \(f\) and has
        \(b\) in its range.
\end{enumerate}
Then, for every \(f_0\in\mathcal F\), there is an isomorphism
\(h:\mathcal M\to\mathcal N\) with \(f_0\subseteq h\), extension taken as
sets of ordered pairs; in particular \(\mathcal M\cong\mathcal N\).
\textit{Status: complete proof.}
\end{theorem}

\begin{proof}
Enumerate \(M=\{a_0,a_1,\ldots\}\) and
\(N=\{b_0,b_1,\ldots\}\), and let \(f_0\) be the given member of
\(\mathcal F\).
Recursively use clause 1 at stage \(2n\) to put \(a_n\) into the domain and
clause 2 at stage \(2n+1\) to put \(b_n\) into the range, obtaining an
increasing sequence \(f_0\subseteq f_1\subseteq f_2\subseteq\cdots\) in
\(\mathcal F\).  The union \(h=\bigcup_{k\in\N}f_k\) of the
resulting increasing sequence is a function because the maps extend one
another, and it is injective because every stage is.  The alternating
requirements make it total and surjective.  Every atomic instance involves
only finitely many arguments and is therefore already preserved and reflected
at some stage.  The union is an isomorphism.  Since
\(f_0\subseteq f_1\subseteq\cdots\) and \(h=\bigcup_{k\in\N}f_k\), the
isomorphism extends \(f_0\).
\end{proof}

See \cite{Hodges1993,Marker2002} for further back-and-forth applications.

\begin{corollary}[Countable saturated structures]
Assume the Axiom of Choice in the ambient set theory.
If \(\mathcal M\) is countable and \(\aleph_0\)-saturated, every finite
elementary map from \(\mathcal M\) to itself extends to an automorphism of
\(\mathcal M\).
\textit{Status: complete proof.}
\end{corollary}

\begin{proof}
Let \(f\) be the given finite elementary map, and let \(\mathcal F\) be
the set of finite elementary maps from \(\mathcal M\) to itself that
extend \(f\).  The set is nonempty, because \(f\in\mathcal F\).  Every
member is a finite partial isomorphism from \(\mathcal M\) to
\(\mathcal M\): elementary maps preserve and reflect all atomic formulas,
and they reflect equality, hence are injective.  Clause 1 of
\cref{thm:back-and-forth-test} holds within \(\mathcal F\): given
\(g\in\mathcal F\) and \(a\in M\), \cref{thm:saturation-homogeneity}
extends \(g\) to an elementary map whose domain also contains \(a\), and
this extension still extends \(f\).  Clause 2 holds by applying the same
theorem to \(g^{-1}\), which is elementary because the inverse of a
finite elementary self-map is elementary: extend \(g^{-1}\) to an
elementary map whose domain contains the prescribed element \(b\) and
invert the result, obtaining a member of \(\mathcal F\) that extends
\(g\supseteq f\) and has \(b\) in its range.
\Cref{thm:back-and-forth-test}, applied with \(f_0\defeq f\), yields an
isomorphism \(h:\mathcal M\to\mathcal M\) with \(f\subseteq h\): an
automorphism extending \(f\).
\end{proof}

\begin{theorem}[Countable categoricity of dense linear order]
\label{thm:dlo-countable-categorical}
Any two countable dense linear orders without endpoints are isomorphic.
Consequently, \(\mathsf{DLO}\) is \(\aleph_0\)-categorical.
\textit{Status: complete proof.}
\end{theorem}

\begin{proof}
Let \(\mathcal M\) and \(\mathcal N\) be countable models, with enumerations
\(M=\{m_0,m_1,\ldots\}\) and \(N=\{n_0,n_1,\ldots\}\).  Construct a chain of
finite order-preserving bijections between finite subsets, starting with
\(f_0=\varnothing\):
\[
  f_0\subseteq f_1\subseteq f_2\subseteq\cdots.
\]
For the extension step, let \(g:A\to B\) be the current finite
order-preserving bijection, where \(A\subseteq M\) and \(B\subseteq N\),
and choose \(a\in M\setminus A\).  Define the finite subsets of \(N\)
\[
 L=\{g(x):x\in A,\ x<^{\mathcal M}a\},\qquad
 U=\{g(x):x\in A,\ a<^{\mathcal M}x\}.
\]
Linearity and \(a\notin A\) give \(L\cap U=\varnothing\) and
\(L\cup U=B\).  Every element of \(L\) is below every element of
\(U\), because \(g\) preserves order.  The following four cases supply
an image \(b\) in \(N\).
\begin{enumerate}
  \item If \(L\ne\varnothing\) and \(U\ne\varnothing\), their
        finiteness gives \(\ell=\max L\) and \(u=\min U\), with
        \(\ell<^{\mathcal N}u\).  Density gives
        \(\ell<^{\mathcal N}b<^{\mathcal N}u\).
  \item If \(L\ne\varnothing\) and \(U=\varnothing\), let
        \(\ell=\max L\).  Since \(\mathcal N\) has no greatest
        element, choose \(b\) with \(\ell<^{\mathcal N}b\).
  \item If \(L=\varnothing\) and \(U\ne\varnothing\), let
        \(u=\min U\).  Since \(\mathcal N\) has no least
        element, choose \(b\) with \(b<^{\mathcal N}u\).
  \item If \(L=U=\varnothing\), then \(B=\varnothing\).
        Choose any \(b\in N\), using the nonempty-domain convention.
\end{enumerate}
In each case \(b\notin B\): it is strictly above every element of
\(L\) and strictly below every element of \(U\), and those sets
partition \(B\).  Therefore \(g\cup\{(a,b)\}\) is an injective
extension.  For every old domain element \(x\), its position below or
above \(a\) agrees exactly with the position of \(g(x)\) below or
above \(b\).  Old pairs retain their order relations, so the extension
preserves and reflects order and is a bijection onto \(B\cup\{b\}\).

At an even stage apply this construction to the first \(m_i\) outside
the current domain.  At an odd stage apply the same four cases to the
inverse current map and the first \(n_i\) outside its domain, using
density and absence of endpoints in \(\mathcal M\); then invert the
extension.  This adds the required \(n_i\) to the range while retaining
every previously assigned pair.  This is the back-and-forth method.

The union \(f=\bigcup_{n\in\N}f_n\) is a function and is injective
because the finite maps form an increasing chain.  Choosing the first
unused enumerated element ensures that every \(m_i\) eventually belongs
to its domain and every \(n_i\) eventually belongs to its range.  Any
two domain elements occur together at some finite stage; their order is
preserved and reflected there and hence by \(f\).  Thus \(f\) is an
isomorphism.  Finally, \((\Q,<)\) is a countably infinite model of
\(\mathsf{DLO}\): rationals are enumerated by integer/positive-integer
pairs, the arithmetic mean lies strictly between two distinct rationals,
and \(q-1<q<q+1\) supplies both endpoint conditions.  Existence at
\(\aleph_0\), together with the isomorphism just constructed, proves
the stated categoricity.
\end{proof}

\begin{hbtheorem}{thm:p03-acfp-uncountable-categorical}{cited=Steinitz1910}{Uncountable categoricity of algebraically closed fields}
Assume the Axiom of Choice in the ambient set theory, and fix \(p=0\) or
a prime.  For every uncountable cardinal \(\kappa\), there is an
algebraically closed field of characteristic \(p\) and cardinality
\(\kappa\), and any two such fields are isomorphic.  Thus
\(\mathsf{ACF}_p\) is \(\kappa\)-categorical for every uncountable
\(\kappa\).  It is not \(\aleph_0\)-categorical.
\end{hbtheorem}

The \emph{prime field} is the smallest subfield; it is isomorphic to
\(\Q\) in characteristic zero and to \(\mathbb Z/p\mathbb Z\) in
characteristic \(p>0\).  A \emph{transcendence basis} over that prime
field is a subset of the given field, algebraically independent over the
prime field, such that the given field is algebraic over the subfield they
generate.  Its cardinality is the \emph{absolute transcendence degree}.
The controlling result is Steinitz's classification by characteristic and
absolute transcendence degree, together with his cardinality calculation
\cite[\S23(6), pp.~300--301, and \S24, p.~301]{Steinitz1910}.
It says that an algebraically closed field is
determined up to isomorphism by these two invariants.  His cardinality
calculation makes the absolute transcendence degree equal to \(\kappa\)
for every uncountable field of cardinality \(\kappa\).  A bijection
between bases, together with the characteristic-preserving isomorphism of
prime fields, induces an isomorphism of their rational-function fields,
and uniqueness of algebraic closures extends it to the closed fields
\cite[\S21(9), p.~287, and \S22(5), p.~289]{Steinitz1910}.

Existence at each uncountable \(\kappa\) uses the source's construction
of a rational-function field over the prime field in \(\kappa\)
independent indeterminates, followed by an algebraic closure.  The same
cardinality calculation gives that closure size \(\kappa\).
Finite or countably infinite transcendence bases instead give countable
algebraically closed fields, and different degrees give nonisomorphic
ones.  In particular, algebraic closures of the prime field and of its
one-indeterminate rational-function field have degrees zero and one.
They are countably infinite: there are only countably many polynomials
over a countable field, each nonzero polynomial has finitely many roots,
and an algebraically closed field cannot be finite.  This explains the
countable-model qualification.  The field classification, existence and
closure theorems are the cited dependencies; their full proofs are not
reproduced in this explanatory application
\cite[\S\S21--24, especially pp.~287--289 and 300--301]{Steinitz1910}.

\begin{theorem}[Ryll--Nardzewski]
Let \(T\) be a deductively complete theory in a countable signature with an infinite
model.  The following are equivalent:
\begin{enumerate}
  \item \(T\) is \(\aleph_0\)-categorical;
  \item for every \(n\geq1\), \(T\) has only finitely many complete
        \(n\)-types over \(\varnothing\); and
  \item for every \(n\geq1\), there are, up to equivalence modulo
        \(T\), only finitely many formulas whose free variables are among
        \(x_1,\ldots,x_n\).
\end{enumerate}
\textit{Status: cited theorem.}
\end{theorem}

See \cite{Hodges1993} for a complete proof.

\begin{theorem}[Morley's categoricity theorem]
Let \(T\) be a deductively complete theory in a countable signature.  If \(T\) is
categorical in one uncountable cardinal, then \(T\) is categorical in every
uncountable cardinal.
\textit{Status: cited theorem.}
\end{theorem}

See \cite{Marker2002} for a complete proof.

\begin{warning}[Completeness versus categoricity]
A deductively complete theory may have many nonisomorphic models of the same
cardinality.  By \cref{def:theory-completeness}, its equivalent semantic
condition says that its models agree on sentences; categoricity says that
models of a specified size are isomorphic.
\end{warning}
